\section{Axis catalogue and column glossary}
\label{app:axes}

\subsection{Prompt templates}

The two executed prompt levels are reproduced verbatim below. Level 1 is also reproduced
because the 360 injected single-agent baseline cells use it; level 2 was removed by the
trim and never ran post-trim.

\subsubsection*{Level 0}

\begin{promptquote}
You are a helpful assistant. Answer the question.
\end{promptquote}

\subsubsection*{Level 1 (single-agent baseline cells only)}

\begin{promptquote}
You are an expert problem solver with broad knowledge across science, mathematics, and
general reasoning. Read the question carefully and provide the most accurate answer you
can. If the question is multiple choice, choose the single best option.
\end{promptquote}

\subsubsection*{Level 3}

\begin{promptquote}
You are a world-class expert problem solver with doctoral-level mastery of the natural
sciences, mathematics, logic, and general reasoning. Your goal is to produce the most
accurate possible answer while reasoning transparently and calibrating your confidence.
\medskip

Follow this strategy on every question:
\medskip

1. COMPREHEND. Restate the question in your own words. Identify the domain, the precise
quantity or fact being asked for, and any traps or ambiguities in the phrasing.
\smallskip

2. RETRIEVE. Surface the relevant principles, theorems, definitions, or empirical facts.
Note any that are commonly misremembered.
\smallskip

3. REASON. Work through the problem step by step. Make each inferential step explicit and
verify it before moving on. For quantitative problems, carry units and sanity-check
magnitudes.
\smallskip

4. CONSIDER ALTERNATIVES. For multiple-choice questions, evaluate every option. Explain
why each distractor is wrong. Watch for options that are true statements but do not
answer the question asked.
\smallskip

5. CHECK. Look for arithmetic slips, sign errors, unit mismatches, and logical gaps. Ask
whether your answer would survive a knowledgeable critic.
\smallskip

6. CALIBRATE. Judge how confident you genuinely are given the strength of your reasoning,
and report that confidence honestly rather than defaulting to high certainty.
\medskip

FORMAT. Show your reasoning, then give your final answer on its own line. For
multiple-choice questions, the final line must name exactly one option letter. Then state
your confidence as a percentage on a separate line.
\medskip

Worked example (format only):\\
Question: Which gas is most abundant in Earth's atmosphere? A. Oxygen B. Nitrogen
C. Carbon dioxide D. Argon\\
Reasoning: Dry air is \textasciitilde78\% N2 by volume, \textasciitilde21\% O2, <1\% Ar,
\textasciitilde0.04\% CO2. The most abundant is nitrogen.\\
Final answer: B\\
Confidence: 99\%
\end{promptquote}

\subsection{Question rendering}

Multiple-choice questions are rendered as the stem, a blank line, the lettered options one
per line, a blank line, and the instruction \code{Respond with the letter of the best
option.} Free-form questions are rendered as the stem followed by \code{Give your final
answer inside \textbackslash boxed\{\}.} Both then append:

\begin{promptquote}
After your answer, on a new line write 'Confidence: X\%' where X is your confidence
(0-100) that your answer is correct.
\end{promptquote}

When thinking is disabled, the user message is additionally prefixed with \code{Think step
by step. Show your reasoning, then give your final answer.} When thinking is enabled this
hint is omitted, because the native thinking mode serves the same purpose.

The deterministic option probe uses a separate completion-style rendering that ends in
\code{Answer: } and is scored at temperature zero with a single generated token.

\subsection{Peer-context rendering}

A rendered peer block is introduced by:

\begin{promptquote}
Here is what other agents concluded. Consider their views, then give your own answer.
\end{promptquote}

Each block then contains \code{Agent \{id\} answered: \{choice\}}, optionally followed by
\code{(stated confidence \{X\}\%)}, and at the plus-CoT level a \code{Reasoning:} field
carrying up to 4{,}000 characters of the peer's chain of thought. Blocks are separated by a
blank line, and the whole rendered context is capped at 4{,}000 exact tokens under the
serving-profile tokenizer. Truncation appends a marker recording the original token count
and a digest of the untruncated text.

\subsection{Column glossary}

Table~\ref{tab:glossary} defines every column of the analysis population.

\intable{tab_glossary}
